%%%%%%%%%%%%%%%%%%%%%%%%%%%%%%%%%%%%%%%%%%%%%%%%%%%%%%%%%%%%%%%%%%%%%%%%%%%%%%%
\documentclass[hyperref={pdfpagelabels=false},compress,table]{beamer} % 在Mac下无法编译
% \documentclass[compress,table]{beamer} % 在Mac下使用
% package for font
\usepackage{fontspec}
\defaultfontfeatures{Mapping=tex-text}  %%如果没有它，会有一些 tex 特殊字符无法正常使用，比如连字符。
\usepackage{xunicode,xltxtra}
\usepackage[BoldFont,SlantFont,CJKnumber,CJKchecksingle]{xeCJK}  % \CJKnumber{12345}: 一万二千三百四十五
\usepackage{CJKfntef}  %%实现对汉字加点、下划线等。
\usepackage{pifont}  % \ding{}
% package for math
\usepackage{amsfonts}

% package for graphics
\usepackage[americaninductors,europeanresistors]{circuitikz}
\usepackage{tikz}
\usetikzlibrary{plotmarks}  % placements=positioning
\usepackage{graphicx}  % \includegraphics[]{}
\usepackage{subfigure}  %%图形或表格并排排列
% package for table
\usepackage{colortbl,dcolumn}  %% 彩色表格
\usepackage{multirow}
\usepackage{multicol}
\usepackage{booktabs}
% package for code
\usepackage{fancyvrb}
\usepackage{listings}

% \usepackage{animate}
% \usepackage{movie15}

%%%%%
% setting for beamer
\usetheme{default} % Madrid（常用）, Copenhagen, AnnArbor, boxes（白色）, Frankfurt，Berkeley
\useoutertheme[subsection=true]{miniframes} % 使用Berkeley时注释本行
\usecolortheme{sidebartab}
\usefonttheme{serif}  %%英文使用衬线字体
% \setbeamertemplate{background canvas}[vertical
% shading][bottom=white,top=structure.fg!7] %%背景色，上25%的蓝，过渡到下白。
\setbeamertemplate{theorems}[numbered]
\setbeamertemplate{navigation symbols}{}  %% 去掉页面下方默认的导航条
\setbeamercovered{transparent}  %设置 beamer 覆盖效果

% 设置标题title背景色
% \setbeamercolor{title}{fg=black, bg=lightgray!60!white}
\setbeamercolor{title}{fg=white, bg=black!70!white}

% 设置每页小LOGO
\pgfdeclareimage[width=1cm]{ouc}{figures/static/ouc.pdf}
\logo{\pgfuseimage{ouc}{\vspace{-20pt}}}

% setting for font
%%\setCJKmainfont{Adobe Kaiti Std}
\setCJKmainfont{SimSun} 
%% \setCJKmainfont{FangSong_GB2312} 
%% \setmainfont{Apple Garamond}  %%苹果字体没有SmallCaps
\setCJKmainfont{SimSun} 
%FUNNY%\setCJKmainfont{DFPShaoNvW5-GB}  %%华康少女文字W5(P)
%FUNNY%\setCJKmainfont{FZJingLeiS-R-GB}  %%方正静蕾体
%FUNNY%\setmainfont{Purisa}
%\setsansfont[Mapping=tex-text]{Adobe Song Std}
     %如果装了Adobe Acrobat，可在font.conf中配置Adobe字体的路径以使用其中文字体。
     %也可直接使用系统中的中文字体如SimSun、SimHei、微软雅黑等。
     %原来beamer用的字体是sans family；注意Mapping的大小写，不能写错。
     %设置字体时也可以直接用字体名，以下三种方式等同：
     %\setromanfont[BoldFont={黑体}]{宋体}
     %\setromanfont[BoldFont={SimHei}]{SimSun}
     %\setromanfont[BoldFont={"[simhei.ttf]"}]{"[simsun.ttc]"}
% setting for graphics
\graphicspath{{figures/}}  %%图片路径
\renewcommand\figurename{图}

% setting for pdf
\hypersetup{% pdfpagemode=FullScreen,%
            pdfauthor={Xiaodong Wang},%
            pdftitle={Title},%
            CJKbookmarks=true,%
            bookmarksnumbered=true,%
            bookmarksopen=false,%
            plainpages=false,%
            colorlinks=true,%
            citecolor=green,%
            filecolor=magenta,%
            linkcolor=blue,%red(default)
            urlcolor=cyan}

% setting for fontspec
\XeTeXlinebreaklocale "zh"  %%表示用中文的断行
\XeTeXlinebreakskip = 0pt plus 1pt minus 0.1pt  %%多一点调整的空间
%%%%%

% font setting by xeCJK
\setCJKfamilyfont{NSimSun}{NSimSun}
\newcommand{\song}{\CJKfamily{NSimSun}}
%%%\setCJKfamilyfont{AdobeSongStd}{Adobe Song Std}
%%%\newcommand{\AdobeSong}{\CJKfamily{AdobeSongStd}}
\setCJKfamilyfont{FangSong}{FangSong_GB2312}
\newcommand{\fang}{\CJKfamily{FangSong}}
%%%\setCJKfamilyfont{AdobeFangsongStd}{Adobe Fangsong Std}
%%%\newcommand{\AdobeFang}{\CJKfamily{AdobeFangsongStd}}
\setCJKfamilyfont{SimHei}{SimHei}
\newcommand{\hei}{\CJKfamily{SimHei}}
%%%\setCJKfamilyfont{AdobeHeitiStd}{Adobe Heiti Std}
%%%\newcommand{\AdobeHei}{\CJKfamily{AdobeHeitiStd}}
\setCJKfamilyfont{KaiTi}{KaiTi}
\newcommand{\kai}{\CJKfamily{KaiTi}}
%%%\setCJKfamilyfont{AdobeKaitiStd}{Adobe Kaiti Std}
\newcommand{\AdobeKai}{\CJKfamily{AdobeKaitiStd}}
\setCJKfamilyfont{LiSu}{LiSu}
\newcommand{\li}{\CJKfamily{LiSu}}
\setCJKfamilyfont{YouYuan}{YouYuan}
\newcommand{\you}{\CJKfamily{YouYuan}}
\setCJKfamilyfont{FZJingLei}{FZJingLeiS-R-GB}
\newcommand{\jinglei}{\CJKfamily{FZJingLei}}
\setCJKfamilyfont{MSYH}{Microsoft YaHei}
\newcommand{\msyh}{\CJKfamily{MSYH}}

% 自定义颜色
\def\Red{\color{red}}
\def\Green{\color{green}}
\def\Blue{\color{blue}}
\def\Mage{\color{magenta}}
\def\Cyan{\color{cyan}}
\def\Brown{\color{brown}}
\def\White{\color{white}}
\def\Black{\color{black}}

\lstnewenvironment{xmlCode}[1][]{% for Java
  \lstset{
    basicstyle=\tiny\ttfamily,%
    columns=flexible,%
    framexleftmargin=.7mm, %
    % frame=shadowbox,%
    % rulesepcolor=\color{cyan},%
     frame=single,%
    backgroundcolor=\color{white},%
    xleftmargin=4\fboxsep,%
    xrightmargin=4\fboxsep,%
    numbers=left,numberstyle=\tiny,%
    numberblanklines=false,numbersep=7pt,%
    language=xml, %
    }\lstset{#1}}{}

\lstnewenvironment{javaCode}[1][]{% for Java
  \lstset{
    basicstyle=\tiny\ttfamily,%
    columns=flexible,%
    framexleftmargin=.7mm, %
    frame=shadowbox,%
    rulesepcolor=\color{cyan},%
    % frame=single,%
    backgroundcolor=\color{white},%
    xleftmargin=4\fboxsep,%
    xrightmargin=4\fboxsep,%
    numbers=left,numberstyle=\tiny,%
    numberblanklines=false,numbersep=7pt,%
    language=Java, %
    }\lstset{#1}}{}

\lstnewenvironment{shCode}[1][]{% for Java
  \lstset{
    basicstyle=\scriptsize\ttfamily,%
    columns=flexible,%
    framexleftmargin=.7mm, %
    frame=shadowbox,%
    rulesepcolor=\color{brown},%
    % frame=single,%
    backgroundcolor=\color{white},%
    xleftmargin=4\fboxsep,%
    xrightmargin=4\fboxsep,%
    numbers=left,numberstyle=\tiny,%
    numberblanklines=false,numbersep=7pt,%
    language=sh, %
    }\lstset{#1}}{}

\newcommand\ask[1]{\vskip 4bp \tikz \node[rectangle,rounded corners,minimum size=6mm,
  fill=white,]{\Cyan \includegraphics[height=1.5cm]{question} \Large \msyh #1};}

\newcommand\wxd[1]{\vskip 4bp \tikz \node[rectangle,minimum size=6mm,
  fill=blue!60!white,]{\White \ding{118} \msyh #1};}

\newcommand\xyy[1]{\vskip 2bp \tikz \node[rectangle,minimum size=3mm,
  fill=black!80!white,]{\White \msyh\scriptsize #1};}

\newcommand\cxf[1]{\vskip 4bp \tikz \node[rectangle,rounded corners,minimum size=6mm,
  fill=orange!60!white,]{\White \ding{42} \msyh #1};}

\newcommand\samp[1]{\vskip 2bp \tikz \node[rectangle,minimum size=3mm,
  fill=white!100!white,]{\Mage\msyh \small CODE \ding{231} \Black #1};\vskip -8bp}

\newcommand\zhyfly[1]{\tikz \node[rectangle,rounded corners,minimum size=6mm,ball color=red!25!blue,text=white,]{#1};}

\setbeamerfont{frametitle}{series=\msyh} % 修改Beamer标题字体

\makeatletter
\newcommand{\Extend}[5]{\ext@arrow 0099{\arrowfill@#1#2#3}{#4}{#5}}
\makeatother


%%%%%%%%%%%%%%%%%%%%%%%%%%%%%%%%%%%%%%%%%%%%%%%%%%%%%%%%%%%%%%%%%%%%%%%%%%%%%%%
% \titlepage
\title[KevinW@OUC]{\hei {\huge Java EE企业应用系统设计}\\  
HTTP请求处理编程}
\author[王晓东]{王晓东\\
  \href{mailto:wangxiaodong@ouc.edu.cn}{\footnotesize wangxiaodong@ouc.edu.cn}}
\institute[中国海洋大学]{\small 中国海洋大学}
\date{\today}
\titlegraphic{\vspace{-6em}\includegraphics[height=6cm]{static/ouc.pdf}\vspace{-6em}}
%%%%%
\begin{document}
%% Delete this, if you do not want the table of contents to pop up at
%% the beginning of each subsection:
\AtBeginSection[]{                              % 在每个Section前都会加入的Frame
  \frame<handout:0>{
    \frametitle{\textbf{\hei 接下来…}}
    \tableofcontents[currentsection]
  }
}  %

\AtBeginSubsection[]                            % 在每个子段落之前
{
  \frame<handout:0>                             % handout:0 表示只在手稿中出现
  {
    \frametitle{\textit{\hei 接下来…}}\small
    \tableofcontents[current,currentsubsection] % 显示在目录中加亮的当前章节
  }
}
 \frame{\titlepage}
 %%%%%%%%%%%%%%%%%%%%%%%%%%%%%%%%%%%%%%%%%%%%%%%%
\begin{frame}
  \frametitle{学习目标}

  \begin{enumerate}
  \item 掌握HTTP协议的特点以及HTTP请求中包含哪些信息。
  \item 理解Java HTTP请求对象的类型及其生命周期，掌握请求对象的功能。
  \item 学习并实践掌握部分请求对象方法的用法。
  \end{enumerate}  
\end{frame}
%%%%%%%%%%%%%%%%%%%%%%%%%%%%%%%%%%%%%%%%%%%%%%%%%%%%%%%%%%%%%%%%%%%%%%%%%%%%%%%
 
\section*{大纲}
\frame{\frametitle{大纲} \tableofcontents}

\section{HTTP请求内容}

\begin{frame}[fragile] % [fragile]参数使得能够插入代码
  \frametitle{Web工作模式}

  Web通常使用{\bf\Red 请求—响应}模式。

  \begin{itemize}\kai
  \item 客户端（浏览器）向服务器发出HTTP请求，在HTTP请求中包含传递到服务器的数据；
  \item Web服务器接收到请求，对请求进行处理。
  \item Web服务器使用HTTP向客户端发送响应。
  \item 客户端接收到响应后，进行显示或页面跳转。
  \end{itemize}
\end{frame}

\begin{frame}[fragile] % [fragile]参数使得能够插入代码
  \frametitle{HTTP请求中包含的信息} 

  HTTP请求中包含的信息包括两部分：{\kai\Red 请求头和请求体}。

  \tta{请求头}

  \begin{xmlCode}
    GET /articles/news/today.jsp HTTP/1.1 
    Accept: */*
    Accept-Language: en-us 
    Connection: Keep-Alive 
    Host: localhost
    Referer: http://localhost/links.asp
    User-Agent: Mozilla/4.0 (compatible; MSIE 5.5; Windows NT 5.0)
    Accept-Encoding:gzip, deflate
    ...
  \end{xmlCode}
\end{frame}

\begin{frame}[fragile] % [fragile]参数使得能够插入代码
  \frametitle{HTTP请求中包含的信息} 

  \tta{HTTP请求头标记和说明}
  
  \begin{description}\kai
  \item[User-Agent] 浏览器的机器环境
  \item[Accept] 浏览器支持哪些MIME数据类型
  \item[Accept-Charset] 浏览器支持的字符编码
  \item[Accept-Encoding] 浏览器支持哪种数据压缩格式
  \item[Accept-Language] 浏览器指定的语言环境
  \item[Host] 浏览器访问的主机名
  \item[Referer] 浏览器是从哪个页面来的
  \item[Cookie] 浏览器保存的cookie对象
  \end{description}

  {\kai\Red Java EE Web组件Servlet和JSP中可以使用请求对象的方法读取这些
    请求内容，进而进行相应的处理。}
\end{frame}

\begin{frame}[fragile,fragile] % [fragile]参数使得能够插入代码
  \frametitle{HTTP请求中包含的信息} 

  \tta{请求体} 

  每次HTTP请求时，在请求头之后会有一个{\bf\Red 空行}，接下来是请求中包含的提交数据，即{\bf\Blue 请求体}。

\end{frame}

\begin{frame}[fragile,fragile] % [fragile]参数使得能够插入代码
  \frametitle{HTTP请求中包含的信息} 

  \ttc{\ding{182} GET请求}

  无请求体，请求数据直接在请求的URL地址中，作为URL的一部分发送给Web服务器。
  
  \begin{xmlCode}
    http://localhost:8080/webapp/login.do?id=9001&pass=9001
  \end{xmlCode}

  \pause
  
  \begin{itemize}
  \item 请求体为空，提交数据直接在URL上，作为请求头部分传输到Web服务器，
    通过URL的QueryString部分能得到提交的参数数据。
  \item 此种方式对提交数据的大小有限制，不同浏览器会有所不同，
    如IE为2083字节。GET请求时数据会出现在URL中，保密性差，实际编程中要
    尽量避免。
  \end{itemize}

\end{frame}

\begin{frame}[fragile,fragile] % [fragile]参数使得能够插入代码
  \frametitle{HTTP请求中包含的信息} 

  \ttc{\ding{183} POST请求}

  \begin{itemize}
  \item 请求体数据单独打包为数据块，通过Socket直接传递到Web服务器端，数
    据不会在地址栏出现。
  \item 可以提交大的数据，包括二进制文件，实现文件上传功能。
  \item 原则上POST请求对提交的数据没有大小限制。
  \end{itemize}
\end{frame}

\section{Java EE请求对象}

\begin{frame}
  \frametitle{请求对象类型与生命周期} 
  \tta{请求对象接口类型}
  \begin{itemize}
  \item Java EE规范中的通用请求对象要实现接口javax.servlet.ServeltRequest
  \item HTTP请求对象要实现接口\\{\bf\Red javax.servlet.http.HttpServletRequest}
  \end{itemize}

  \tta{请求对象生命周期}

  在Java Web组件开发中，不需要Servlet或JSP自己创建请求对象，它们
  由{\bf\Blue Web容器自动创建}，并传递给Servlet和JSP的服务方
  法doGet和doPost，在服务处理方法中直接使用请求对象即可。
\end{frame}

\begin{frame}
  \frametitle{请求对象类型与生命周期} 
  \tta{请求对象创建}

  每次Web服务器接收到HTTP请求时，会自动创建实现HttpServletRequest接口的
  对象。在创建该对象之后，Web服务器将请求头和请求体信息写入请求对
  象，Servlet和JSP可以通过请求对象的方法取得这些请求信息，继而可以取得
  用户提交的数据。

  \tta{请求对象销毁}

  当Web服务器处理HTTP请求，向客户端发送HTTP响应结束后，会自动销毁请求对
  象，保存在请求对象中的数据随即丢失。当下次请求时新的请求对象又会被创
  建。
\end{frame}

\begin{frame}[fragile] % [fragile]参数使得能够插入代码
  \frametitle{请求对象功能与方法} 

  \tta{请求对象方法一般分类}
  
  \begin{itemize}
  \item 取得请求头信息；
  \item 取得请求体中包含的提交参数数据，包含表单元素或地址栏URL的参数；
  \item 取得客户端的有关信息，如请求协议、IP地址和端口等；
  \item 取得服务器端的相关信息，如服务器的IP等；
  \item 取得请求对象的属性信息，用于在一个请求的转发对象之间传递数据。
  \end{itemize}
\end{frame}

\begin{frame}[fragile] % [fragile]参数使得能够插入代码
  \frametitle{取得请求头} 

  \begin{itemize}
  \item String getHeader(String name)
    \begin{javaCode}
      String browser = request.getHeader("User-Agent");
    \end{javaCode}

  \item int getIntHeader(String name)
    \begin{javaCode}
      int size = request.getIntHeader("Content-Length");
    \end{javaCode}

  \item long getDateHeader(String name)
    \begin{javaCode}
      long datetime = request.getDateHeader("If-Modified-Since");
    \end{javaCode}

  \item Enumeration getHeaderNames()
    
    \begin{javaCode}
      for (Enumeration e = request.getHeaderNames(); e.hasMoreElements(); ) {
        String headerName = (String) e.nextElement();
        System.out.println("Name = " + headerName);
      }
    \end{javaCode}
  \end{itemize}
\end{frame}

\begin{frame}[fragile,fragile] % [fragile]参数使得能够插入代码
  \frametitle{取得请求中包含的提交参数数据} 

  \begin{itemize}
  \item 在Web开发中，用户通过表单或URL参数将客户端数据提交到服务器端，
    这些数据被Web服务器自动封装到请求对象中。
  \item Web组件Servlet和JSP可以通过请求对象获得用户提交的数据。
  \end{itemize}
\end{frame}

\begin{frame}[fragile,fragile] % [fragile]参数使得能够插入代码
  \frametitle{取得请求中包含的提交参数数据} 
  
  \tta{String getParameter(String name)} 取得表单或URL参数中指定名称的数据值。

  \ttc{表单}

  \begin{xmlCode}
    Product Name: <input type="text" name="productName" />
  \end{xmlCode}

  \ttc{URL参数}
  
  \begin{xmlCode}
    productSearch.do?productName=Acer    
  \end{xmlCode}

  以下代码可以取得以上的参数名为productName的数据：
  
  \begin{javaCode}
    String productName = request.getParameter("productName");
  \end{javaCode}
\end{frame}

\begin{frame}[fragile,fragile] % [fragile]参数使得能够插入代码
  \frametitle{取得请求中包含的提交参数数据} 

  \tta{String[] getParameterValues(String name)} 取得指定名称的参数数据数组，如复选框和复选列表。

  \ttc{表单}
  
  \begin{xmlCode}
    爱好：
    <input type="checkbox" name="behave" value="旅游">旅游
    <input type="checkbox" name="behave" value="读书">读书
    <input type="checkbox" name="behave" value="体育">体育
  \end{xmlCode}

  如下代码取得选定的爱好：

  \begin{javaCode}
    String[] behaves = request.getParameterValues("behave");
    for (int i = 0; i < behaves.length; i++) {
      out.println(behaves[i]);
    }
  \end{javaCode}
\end{frame}

\begin{frame}[fragile] % [fragile]参数使得能够插入代码
  \frametitle{取得请求中包含的提交参数数据} 
  \tta{Enumeration getParameterNames()} 取得所有请求的参数名称。

  \begin{javaCode}
    for (Enumeration enums = request.getParameterNames(); enums.hasMoreElements();) {
      String paramName = (String) enums.nextElement();
      System.out.println(paramName);
    }
  \end{javaCode}
\end{frame}

\begin{frame}[fragile] % [fragile]参数使得能够插入代码
  \frametitle{取得请求中包含的提交参数数据} 

  \tta{Map getParameterMap()} 取得所有请求的参数名和值，包装在一个Map对
  象中，可以使用这个对象同时取得所有的参数名和参数值。

  \begin{javaCode}
    Map params = request.getParameterMap();
    Set names = params.keySet();
    for (Object o: names) {
      String paramName = (String) o;
      out.print(paramName + " = " + params.get(paramName) + "<br/>");
    }
  \end{javaCode}
\end{frame}

\begin{frame}[fragile,fragile] % [fragile]参数使得能够插入代码
  \frametitle{取得请求中包含的提交参数数据} 

  \tta{ServletInputStream getInputStream() throws IOException} 取得客户
  端的输入流。

  \begin{itemize}
  \item 当用户提交的数据中包含文件上传时，提交的数据可以以二进制编码方
    式提交到服务器。
  \item 当表单既有文本字段还有文件上传时，就需要对此二进制流进行解析，
    从而分离出文本和上传文件。
  \item 可以使用第三方框架/Jar包实现上传文件的处理，如{\kai\Blue Apache的Common
      upload组件}。
  \end{itemize}

  \notice{注意}

  {\kai 当使用getParameter()方法后，就无法使用getInputStream()方法，反之亦然。}


\end{frame}

\begin{frame}[fragile,fragile] % [fragile]参数使得能够插入代码
  \frametitle{取得请求中包含的提交参数数据} 

  \begin{xmlCode}
    <form action="" method="" enctype="multipart/form-data">
      姓名：<input type="text" name="name" /><br />
      照片：<input type="file" name="photo" /><br />
      <input type="submit" value="提 交" />
    </form>
\end{xmlCode}
\end{frame}

\begin{frame}[fragile] % [fragile]参数使得能够插入代码
  \frametitle{取得其他客户端信息} 

  \ttc{String getRemoteHost()} 取得请求客户的主机名。

  \ttc{String getRemoteAddr()} 取得请求客户端的IP地址。

  \ttc{int getRemotePort()} 取得请求客户的端口号。

  \ttc{String getProtocol()} 取得请求的协议。
  
  \ttc{String getContentType()} 取得请求体的MIME内容类型。

  \ttc{int getContentLength()} 取得请求体为二进制流时请求体的长度。

  \ttc{String getProtocol()} 取得请求的协议，一般为HTTP。
\end{frame}

\begin{frame}[fragile] % [fragile]参数使得能够插入代码
  \frametitle{取得服务器端信息} 

  \ttc{String getServerName()} 取得服务器的HOST，一般为IP地址。

  \ttc{int getServerPort()} 取得服务器接收端口。
\end{frame}

\begin{frame}
  \frametitle{本节习题}

  \tta{问答题}
  
  \begin{enumerate}
  \item 自行搜索总结HTTP协议的特点，除了GET和POST方法外，HTTP所包含的其
    他方法的含义及用途。
  \item 对HTTP请求头可以使用的标记及其作用进行总结。
  \end{enumerate}

  \tta{小编程}
  \begin{enumerate}
  \item 实践课堂示例代码，完成HTTP GET和POST请求的编程测试，实现服务器
    端对客户端请求参数和表单数据的获取。
  \end{enumerate}
\end{frame}
%%%%%%%%%%%%%%%%%%%%%%%%%%%%%%%%%%%%%%%%%%%%%%%%%%%%%%%%%%

% TKS %%%%%%%%%%%%%%%%%%%%%%%%%%%%%%%%%%%%%%%%%%%%%%
\begin{frame}
\centering
{\Huge \textcolor{blue}{THE END}} \\
\vspace{5mm}
{\Large wangxiaodong@ouc.edu.cn} \\
\end{frame}
%%%%%%%%%%%%%%%%%%%%%%%%%%%%%%%%%%%%%%%%%%%%%%%%%%
\end{document}
